% -------------- EXPERIMENTOS Y ANÁLISIS DE RESULTADOS --------------
% Contenido (sumario de la planificación, §6): E1 a E4, con costos y latencias.
\clearpage
\blank
\tituloapa[Experimentos y análisis de resultados]{EXPERIMENTOS Y ANÁLISIS DE RESULTADOS}
\begingroup\dobleesp
\vspace{6pt}
Este capítulo presenta los resultados de los experimentos en el orden del diseño experimental. En esta versión del documento se presenta solo E1, la medición de la capacidad de detección sobre el framework principal. Sus cifras son preliminares por dos razones fijadas antes del experimento: la auditoría del oráculo sobre cuarenta trazas de E1, que hacen anotadores externos, puede obligar a corregir sus reglas como enmienda y a volver a puntuar el experimento; y la concordancia del evaluador automático con los anotadores, que mide E2, todavía no se conoce.
\par

\subtituloapa{E1: ejecución}
E1 se ejecutó el 18 de octubre de 2026 sobre la edición Claude Agent SDK, con el protocolo congelado en su versión vigente. Se corrieron ocho condiciones, la línea base de referencia, la línea base de control y las seis mutaciones, sobre los veinte escenarios y con cinco repeticiones cada una: 800 conversaciones, cada una con su juicio. La primera condición se corrió sola para confirmar su costo en la consola del proveedor antes de lanzar las demás, que se ejecutaron en dos tandas paralelas. Las conversaciones tomaron unas seis horas y media de reloj, incluida la espera para leer la consola.
\par
No hubo ningún reintento. Ninguna conversación terminó por un fallo del usuario simulado ni por un error del agente, y ningún juicio falló, así que todos los pares de escenario y condición tienen sus cinco repeticiones válidas. Todas las conversaciones terminaron porque el usuario simulado las dio por cerradas, salvo 26 de D6, que se cortaron por el límite de pasos y caen todas en su conjunto expuesto. Trece trazas quedaron fuera del alcance de las reglas del oráculo y el autor las anotó a ciegas, en una sola sesión y en un orden fijado por semilla. Son menos de las que proyectó el piloto.
\par
E1 costó USD 37.89 según la facturación del proveedor: 19.44 del evaluador, 13.26 del agente y 5.19 del modelo menor, que usan el usuario simulado, el agente de D3 y algunas llamadas internas del SDK. El costo del evaluador coincide al centavo con el registro de AgentProof. El del agente es el 40\,\% de lo que estima su SDK, por la caché de prompts, una diferencia que se observó desde el piloto y por la que la tesis toma el costo del agente de la facturación.
\par
\endgroup
